God's family asks to be free from every adversity, yet at the Offertory it
hears of a righteous man whose children, possessions, and bodily health are
taken from him. A servant receives an immense pardon, yet at the end of the
Gospel he is handed over to punishment. These are the two sharp questions
of the Twenty-first Sunday after Pentecost. What does mercy require of the
person who has received it? What does protection mean for people who still
suffer? The Mass neither conceals the wounds nor softens the judgment. It
arms its hearers for a conflict in which the human neighbor is not the
ultimate enemy, and feeds them with what the Postcommunion calls the food
of immortality.

The two readings developed below answer those questions in different ways.
In the first, the king's prior gift establishes the forgiven servant's
obligation to his fellow servant. Chrysostom follows the repeated plea for
patience and the extraordinary difference between what was asked and what
was given; Augustine places every hearer among both debtors and creditors.
In the second, the coexistence of protection and wounds is decisive.
Gregory distinguishes Satan's intention from the power permitted him over
Job; Augustine hears in the persecuted psalmist the voice of Christ's
whole body awaiting salvation. Chrysostom's account of the Christian Exodus
then explains why a people already delivered still needs armor and food.

Mercy and perseverance belong together, but they are not interchangeable
answers. One exposes the danger of using received freedom to oppress
someone else. The other makes it possible to ask for rescue without
concluding that an unanswered cry proves abandonment. The king's servant
is guilty of cruelty; Job is named upright before his affliction. Keeping
those two figures distinct makes both readings more demanding and more
consoling.

\section{The appointed texts}
\begin{maptable}
Introit & \textit{In voluntate tua}; Esth.~13:9--11; Ps.~118:1 & The endangered people appeal to the creator; the psalm verse names their walk in his law.\\
Sunday Collect & \textit{Familiam tuam} & Continuing mercy guards God's family and frees it for good works.\\
Mission Collect & \textit{Deus, qui omnes} & On the mission day, God is asked to send workers and give confidence to their speech.\\
St Luke Collect & \textit{Interveniat pro nobis} & In the admitted commemoration, the Evangelist's cross-bearing accompanies his intercession.\\
Epistle & Eph.~6:10--17 & Divine armor equips resistance to spiritual wickedness, not warfare against flesh and blood.\\
Gradual & Ps.~89:1--2 & God remains refuge through generations, before the mountains and beyond time.\\
Alleluia & Ps.~113:1 & Israel's departure from Egypt recalls a real deliverance and the beginning of a journey.\\
Gospel & Matt.~18:23--35 & A freely pardoned debtor refuses mercy; the king calls him to account.\\
Offertory & \textit{Vir erat}; Job~1--2 & Upright Job suffers losses and bodily wounds under a bounded permission.\\
Sunday Secret & \textit{Suscipe, Domine} & God wills propitiation and the restoration of salvation through the offerings.\\
Mission Secret & \textit{Protector noster} & Christ's redemption for all opens toward a pure offering among the nations.\\
St Luke Secret & \textit{Donis caelestibus} & Heavenly gifts seek free service, healing, and glory.\\
Trinity Preface & \textit{Qui cum unigenito} & Thanksgiving confesses one God, distinct Persons, and equal majesty.\\
Communion & Ps.~118:81,84,86 & Hope in God's word still asks when he will judge persecution and cries for help.\\
Sunday Postcommunion & \textit{Immortalitatis alimoniam} & What the mouth receives must be followed with a pure mind.\\
Mission Postcommunion & \textit{Redemptionis nostrae} & Redemption's nourishment asks for the continuing advance of true faith.\\
St Luke Postcommunion & \textit{Praesta, quaesumus} & The altar's gift is asked to sanctify and safeguard the soul through the saint's intercession.\\
\end{maptable}

For 18 October 2026 the general-calendar Sunday is second class and green.
The mission prayers are added under the Sunday prayers' conclusions on
the penultimate Sunday of October unless the local Ordinary has fixed
another mission day. St Luke's feast yields to the Sunday: his three
commemorative orations follow with separate conclusions at low or
conventual Mass; they are absent at a nonconventual sung Mass. These are
conditions governing the prayers, not a choice of substitute readings.
The Sunday's Gospel and Trinity Preface remain.\footnote{\work{Missale
Romanum} (1962), general rubrics 91,95,107--111,444--446,453,480,483,494,505;
Sunday pp.~413--415; St Luke pp.~689--690; mission prayers [86]--[87].}

\clearpage
